% Метаданные известной на текущий момент версии. Класс psta загружен с
% опцией blind, поэтому сведения об авторе не попадают в reviewer PDF.
% Перед подачей заполните TODO-поля и согласуйте состав авторов.

\Russian
\title[RR-экстраполяция в линейной стохастической аппроксимации]
{Экстраполяция Ричардсона--Ромберга для статистического вывода
в~линейной стохастической аппроксимации}

\author{Гайнанов, {Д}аниил}
\address{\href{https://www.hse.ru}{Национальный исследовательский университет
«Высшая школа экономики»}, Москва, Россия}
% \email{TODO}
% \orcid{TODO}
% \info{TODO: 4--7 строк о статусе, научных интересах и достижениях.}
% \image{pics/TODO-author-photo}
% \CRediT{1,2,3,4,5,6,8,9,10,11}{100\%}

\begin{abstract}
Рассматривается статистический вывод для линейной стохастической аппроксимации
с постоянным шагом и марковским шумом. Усреднение Поляка--Рупперта снижает
дисперсию, но сохраняет смещение порядка шага. Для коррекции центра две
рекурсии с шагами $\alpha$ и $2\alpha$ запускаются на одной траектории и
объединяются экстраполяцией Ричардсона--Ромберга. Основной результат ---
неасимптотическая гауссовская аппроксимация скалярных проекций RR-среднего
после разогрева. При равномерной геометрической эргодичности, ляпуновской
устойчивости, ограниченном шуме и моментной устойчивости случайных матричных
произведений рассмотрен допустимый режим, в котором выполнены зависящие от
порядка момента пороги малости шага. Тогда выбор $\alpha=c/\sqrt n$ и разогрев
порядка $(\alpha a)^{-1}\log^2n$ дают оценку расстояния Колмогорова до
стандартного нормального закона
$C_{\mathrm{mix}}(t_{\mathrm{mix}})
\operatorname{polylog}(n)n^{-1/4}$. Доказательство выделяет ведущий
пуассоновский мартингал, сравнивает ковариационную прокси с асимптотической
долгосрочной ковариацией и соединяет опубликованные $L_p$-оценки отдельных
компонент с локальной леммой переноса после разогрева и сглаживающим
неравенством. В экспериментах на конечных марковских
цепях RR устраняет видимый сдвиг нормированной статистики и возвращает
покрытие скалярных 95\%-х интервалов к уровню 94--95\% без заметного
увеличения ширины. На длинных траекториях оценка долгосрочной дисперсии методом
перекрывающихся блочных средних близка к оракульной.
\end{abstract}
\keywords{линейная стохастическая аппроксимация,
марковский шум, экстраполяция Ричардсона--Ромберга,
гауссовская аппроксимация, оценка Берри--Эссеена,
доверительные интервалы, долгосрочная ковариация}

\English
\title[RR extrapolation in linear stochastic approximation]
{Richardson--Romberg Extrapolation for Statistical Inference
in Linear Stochastic Approximation}

\author{Gainanov, Danil}
\address{\href{https://www.hse.ru/en/}{HSE University}, Moscow, Russia}
% \email{TODO}
% \orcid{TODO}
% \info{TODO: 4--7 lines on status, research interests, and achievements.}

\begin{abstract}
We study statistical inference for constant-stepsize linear stochastic
approximation driven by Markovian observations. Polyak--Ruppert averaging
reduces random fluctuations but retains a stepsize-dependent bias. To correct
the center, we run two recursions with stepsizes $\alpha$ and $2\alpha$ on the
same trajectory and combine their averages by Richardson--Romberg
extrapolation. Our main result is a non-asymptotic Gaussian approximation for
scalar projections of the burned-in RR average. Under uniform geometric
ergodicity, Lyapunov stability, bounded noise, and moment stability of the
random matrix products, consider the admissible regime in which the
moment-dependent stepsize ceilings hold. Then choosing $\alpha=c/\sqrt n$ and
a burn-in of order $(\alpha a)^{-1}\log^2n$ yields a Kolmogorov-distance
bound of order $C_{\mathrm{mix}}(t_{\mathrm{mix}})
\operatorname{polylog}(n)n^{-1/4}$. The proof isolates a leading
Poisson-equation martingale, compares its covariance proxy with the
asymptotic long-run covariance, and combines published componentwise $L_p$
bounds with a burn-in transfer lemma and a smoothing inequality. In finite-state
experiments, extrapolation removes the visible shift of the normalized
statistic and brings scalar 95\% confidence-interval coverage to 94--95\%
without a material increase in width. At long horizons, overlapping batch
means closely reproduce the oracle long-run variance.
\end{abstract}
\keywords{linear stochastic approximation, Markovian noise,
Richardson--Romberg extrapolation, Gaussian approximation,
Berry--Esseen bound, confidence intervals, long-run covariance}

\Russian
